%% ============================================================
%%  Principia Orthogona, Volume Two
%%  Section 3: Equivalence of kappa* and tau
%% ============================================================

\section{Equivalence of the Curvature Threshold and the
  Embodiment Threshold}
\label{sec:threshold-equivalence}

The central bridge theorem of this volume establishes a
quantitative correspondence between two thresholds that appear
independently in the two frameworks:
\begin{itemize}[leftmargin=*]
  \item the \emph{curvature threshold} $\kappa^*$, defined
        geometrically in Volume One as
        $\kappa^*(x) = 1/\operatorname{foc}(x)$, and
  \item the \emph{embodiment threshold} $\tau$, defined
        dynamically in GCM~\cite{Paper1} as
        $\tau = \sqrt{c/\kappa}$ from the stochastic Lyapunov
        condition $\mathcal{L}V \le -cV + \kappa\|\sigma\|^2$.
\end{itemize}

Both thresholds govern a transition from neutral to stable
behavior. The curvature threshold determines when a trajectory
folds and enters the attracting regime. The embodiment threshold
determines when noise ceases to destabilize the orbit and begins
to reinforce it. We show these are the same transition, viewed
from different geometric perspectives.

%% ------------------------------------------------------------
\subsection{Setup and Assumptions}
\label{subsec:setup}

Let $(X,g)$ be a smooth complete Riemannian manifold and let
$\gamma\colon[0,T]\to X$ undergo a generative transition via the
sequence $C\to K\to F\to U$ of Volume One~\cite{PrincipiaI}.
Let $\Gamma\subset X$ be the post-fold invariant set (the
stabilized orbit produced by $U$), and let $V = d_g(\cdot,\Gamma)^2$
be the Lyapunov function of the resulting dm3 system
$\mathfrak{D}=(X,g,Y,\Gamma,V,\sigma)$~\cite{Paper1}.

\begin{assumption}\label{ass:fold-orbit}
The fold operator $F$ produces a unique post-fold orbit $\Gamma$
that is a hyperbolic limit cycle of the post-fold vector field,
with transverse Floquet exponent $\mu_{\max}<0$. The unfolding
operator $U$ selects this orbit via gradient descent on the
Morse stability functional $\Phi$.
\end{assumption}

This assumption is satisfied in all three canonical examples of
Volume One (elastic rod, double-well gradient flow, saddle-node
bifurcation) and in the dm3 toy model of Paper~2~\cite{Paper2}.

%% ------------------------------------------------------------
\subsection{The Threshold Correspondence Theorem}
\label{subsec:main-theorem}

\begin{theorem}[Threshold Correspondence]
\label{thm:threshold}
Under Assumption~\ref{ass:fold-orbit}, let $\kappa^*$ be the
curvature threshold of Volume One and $\tau$ be the embodiment
threshold of~\cite{Paper1}. Then:
\begin{enumerate}[label=\textup{(\roman*)}]
  \item \textbf{Geometric to dynamic:} The condition
        $\kappa(s)\uparrow\kappa^*$ is equivalent to the
        transverse Floquet exponent becoming negative:
        $\mu_{\max}<0$. Both conditions mark the onset of
        transverse contraction toward $\Gamma$.
  \item \textbf{Dynamic to stochastic:} The condition
        $\mu_{\max}<0$ implies the existence of a finite
        embodiment threshold
        \[
          \tau = \sqrt{c/\kappa} > 0,
        \]
        where $c=|\mu_{\max}|$ is the Lyapunov rate and
        $\kappa$ is the noise amplification constant of
        Axiom~5 of~\cite{Paper1}.
  \item \textbf{Quantitative form:} In the contact normal form
        \textup{(\cite{Paper1}, Theorem~C)}, the embodiment
        threshold satisfies
        \[
          \tau(z) = \sqrt{c(z)/\kappa}
          = 2\sqrt{1-e^{-\beta z}},
        \]
        which increases monotonically from $\tau(0)=0$
        (at the fold, neutral stability) to
        $\tau(\infty) = 2\sqrt{c_\infty/\kappa}$
        (asymptotic stability).
\end{enumerate}
\end{theorem}

\begin{proof}
\textit{Claim (i).}
In the contact normal form~\eqref{eq:nf-rho}, the transverse
dynamics are $\dot\rho = \mu_{\max}(1-e^{-\beta z})\rho +
O(\rho^2)$. The condition $\kappa\uparrow\kappa^*$ corresponds
to the fold event at $s_0$: prior to the fold, $\kappa<\kappa^*$
and trajectories are non-injective in the normal direction only
after the fold. In the contact extension, the post-fold regime
is characterized by $z>0$ (action has accumulated), at which
point $\mu_{\max}(1-e^{-\beta z})<0$. Thus the curvature
threshold $\kappa^*$ and the sign change of the transverse
eigenvalue coincide: both mark the onset of injectivity loss
followed by attraction.

\textit{Claim (ii).}
From Axiom~5 of~\cite{Paper1}, the stochastic Lyapunov condition
$\mathcal{L}V \le -cV + \kappa\|\sigma\|^2$ holds with
$c = |\mu_{\max}|(1-e^{-\beta z}) > 0$ for $z>0$. Since
$c>0$ and $\kappa>0$, the embodiment threshold
$\tau = \sqrt{c/\kappa}$ is finite and positive. It exists
if and only if $\mu_{\max}<0$, i.e., if and only if the fold
has been completed.

\textit{Claim (iii).}
In the toy model of Paper~2~\cite{Paper2}, direct computation
gives $c(z)=4(1-e^{-z})$ and $\kappa=1$, so
$\tau(z)=2\sqrt{1-e^{-z}}$. At $z=0$: $\tau=0$ (neutral,
fold has just occurred). As $z\to\infty$: $\tau\to 2$
(asymptotic threshold). This is Proposition~\ref{prop:axioms}
of Paper~2 and Lemma~\ref{lem:transverse} thereof.
\end{proof}

%% ------------------------------------------------------------
\subsection{Interpretation}
\label{subsec:interpretation}

Theorem~\ref{thm:threshold} has three consequences worth
making explicit.

\paragraph{The fold creates the threshold.}
The embodiment threshold $\tau$ does not exist before the fold:
for $\kappa<\kappa^*$, the orbit has not yet formed and there is
no invariant set to be stabilized. The fold event at $\kappa^*$
is precisely the geometric act that creates the conditions for
stochastic stability. In this sense, $\kappa^*$ is the
\emph{geometric precursor} of $\tau$.

\paragraph{The action variable records the fold.}
The contact variable $z$ measures how long the system has been
in the post-fold regime. At $z=0$ (immediately after the fold),
$\tau=0$: the orbit is neutrally stable and any noise will
displace it. As $z$ grows, $\tau$ grows: the orbit becomes
increasingly noise-tolerant. This is the embodiment mechanism
expressed as a theorem: \emph{stability is earned by traversing
the orbit}.

\paragraph{The thresholds are not numerically equal.}
Theorem~\ref{thm:threshold} establishes a structural
correspondence between $\kappa^*$ and $\tau$, not a numerical
identity. In general, $\kappa^*$ is a geometric quantity
(units: inverse length) and $\tau$ is a dynamical quantity
(units: noise amplitude). Their correspondence is a statement
about which physical transition each governs, not about their
numerical values.

%% ------------------------------------------------------------
\subsection{Explicit Verification in the Toy Model}
\label{subsec:toy-verification}

In the dm3 toy model of Paper~2~\cite{Paper2}, the threshold
correspondence is verified explicitly:

\begin{itemize}[leftmargin=*]
  \item The limit cycle $\Gamma=\{r=1\}$ is produced by the
        fold of the radial dynamics at the critical radius
        $r=1$ (the curvature threshold in the radial direction).
  \item The transverse Floquet exponent is $\mu_{\max}=-2$,
        confirmed by linearization
        $\dot\rho=-2(1-e^{-z})\rho+O(\rho^2)$.
  \item The embodiment threshold is
        $\tau(z)=2\sqrt{1-e^{-z}}\to 2$ as $z\to\infty$.
  \item At $z=0$: $\tau=0$, confirming neutral stability
        at the fold point.
  \item The stability radius is
        $\varepsilon_0 = |\mu_{\max}|/2(1+\|\Hess V\|) = 1/3$,
        computed from Theorem~D of~\cite{Paper1}.
\end{itemize}

All values are consistent with Theorem~\ref{thm:threshold}
and with the global attractor analysis of Paper~2 (Theorem~A
thereof).

%% ------------------------------------------------------------
\subsection{Summary}
\label{subsec:section3-summary}

The curvature threshold $\kappa^*$ of Volume One and the
embodiment threshold $\tau$ of GCM are structurally equivalent:
both mark the onset of transverse contraction following a
fold event. The contact action variable $z$ mediates the
transition between them, recording accumulated post-fold
dissipation and increasing $\tau$ monotonically from $0$ at
the fold to its asymptotic value.

This correspondence is the quantitative content of the claim
made in Volume One (\S13, Proposition~10.1) that $\kappa^*$
is the geometric precursor of $\tau$. It is now stated and
proved as Theorem~\ref{thm:threshold}.

The next section uses this correspondence to construct the
full discrete contact realization of the operator sequence
$C\to K\to F\to U$ and verify it against the dm3 global
dynamics.
